\documentclass{article}

\usepackage[a4paper]{geometry}
\usepackage[english]{babel}
\usepackage{parskip}

\usepackage{amsmath, amssymb}
\usepackage{graphicx}
\usepackage{subcaption}
\usepackage{float}
\usepackage[hidelinks]{hyperref}

\usepackage{listings}
\usepackage{xcolor}

\definecolor{codegreen}{rgb}{0,0.6,0}
\definecolor{codegray}{rgb}{0.5,0.5,0.5}
\definecolor{codepurple}{rgb}{0.58,0,0.82}
\definecolor{backcolour}{rgb}{0.95,0.95,0.92}

\lstdefinestyle{mystyle}{
  backgroundcolor=\color{backcolour},
  commentstyle=\color{codegreen},
  keywordstyle=\color{magenta},
  numberstyle=\tiny\color{codegray},
  stringstyle=\color{codepurple},
  basicstyle=\ttfamily\footnotesize,
  breaklines=true,
  captionpos=b,
  keepspaces=true,
  numbers=left,
  numbersep=5pt,
  showstringspaces=false,
  tabsize=2
}
\lstset{style=mystyle}

\newcommand{\CodeDir}{Code}
\newcommand{\PlotDir}{Plots}

\title{NURa Hand-in 1}
\author{Milan Tol (s3162168)}
\date{1st March 2026}

\begin{document}
\maketitle


\section{Poisson Distribution}

% Add the full code below.
The full code used for this question:
\lstinputlisting[language=Python]{\CodeDir/Poisson_code.txt}

% Show your results.
The results of $P_\lambda(k)$ for selected values of $k$ and $\lambda$ are shown in Table~\ref{tab:poisson}.

\begin{table}[H]
\centering
\caption{Poisson probability distribution for selected $k$ and $\lambda$.}
\label{tab:poisson}
\begin{tabular}{|c|c|c|}
\hline
$\lambda$ & $k$ & $P_\lambda(k)$ \\ 
\hline
\input{Poisson_output.txt} \\
\hline
\end{tabular}
\end{table}




\section{Vandermonde Matrix and Interpolation}

\subsection{(2a) LU decomposition}
\begin{figure}[H]
  \centering
  \includegraphics[width=0.80\textwidth]{\PlotDir/vandermonde_sol_2a.pdf}
  \caption{Polynomial fit evaluated using LU decomposition. Top: data points and interpolated curve. Bottom: absolute error at the data points on a log scale.}
  \label{fig:vandermonde_2a}
\end{figure}



\subsection{(2b) Neville's algorithm}
\begin{figure}[H]
  \centering
  \includegraphics[width=0.80\textwidth]{\PlotDir/vandermonde_sol_2b.pdf}
  \caption{Interpolation using Neville's algorithm. Top: data points and interpolated curve. Bottom: absolute error at the data points on a log scale.}
  \label{fig:vandermonde_2b}
\end{figure}



\begin{figure}[H]
  \centering
  \includegraphics[width=0.80\textwidth]{\PlotDir/Comparison_LU_neville.pdf}
  \caption{Comparison between the LU decomposition solution and the Neville interpolator. Top: interpolated curves. 
  Bottom: absolute difference of the two interpolators on a log scale.}
  \label{fig:LU_neville_comparison}
\end{figure}


\subsection{(2c) Improving the LU decomposition}
\begin{figure}[H]
  \centering
  \includegraphics[width=0.80\textwidth]{\PlotDir/vandermonde_sol_2c.pdf}
  \caption{LU-based solution with iterative refinement (showing iterations 0, 1 and 10). Top: interpolated curves. Bottom: absolute error at the data points on a log scale.}
  \label{fig:vandermonde_2c}
\end{figure}




\subsection{(2d) Timing}

Timing results (average per run):
\begin{itemize}
  \input{Execution_times.txt}
\end{itemize}




\subsection{Conclusions}

We can see that for larger x the error increases substantially, almost like a power law. This is because we are evaluating a 19th order polynomial.
The coefficients of the larger powers of x: $x^\alpha$, will contain some level of error due to roundoff. 
When evaluating the polynomial at large x, we scale this error by a factor of $x^\alpha$. 

The Neville interpolator returns the same polynomial fit as the scheme that implemented solving the linear system in (a).
This is not suprising, since Neville interpolation up to 20th order produces a 19th order polynomial, which also passes through all the 20 points.
Since we have 20 constraints and 20 data points, such polynomials are unique. 
However, the difference from the interpolated value compared to the value returned by Neville interpolation is practically 0 up to some roundoff error for 
\textbf{all} values of x. The roundoff error does not get scaled by some factor dependent on x, unlike the LU scheme. 
This makes sense when looking at how Neville interpolation works.
For each point we interpolate the data separately. 
At each data point $(x_i, y_i)$, for each iteration, we have an expression of the form $y = y_i + (x-x_i)*C(x)$,
where $C$ is a function determined by the other data points. This becomes $y_i$ up to some roundoff error.
However, the price of this precision is speed. Since we must do the interpolation for every point separately the algorithm is almost two orders of magnitude slower!

a comparison between the two interpolators in Figure \ref{fig:LU_neville_comparison}, we see that
in between the two data points the two interpolators do not match all, compared to how well the two methods match the data. 
Why is this? Well its because Neville interpolation is technically not plotting a clean 19th order polynomial. 
For each point at which the neville interpolator is evaluated, it is doing a new fit to the data, 
using the point at which it is being evaluated as $x=0$ (that is also why the errors are so small).

Is the precision worth the computation time? The answer is (almost) all cases is no.
In most practical cases, there is some error associated with the data itself. 
Fitting perfectly to noisy data is not physical. 
One exception might be when we would extend this to a 100th+ order polynomial. 
In that case the LU decomposition approach would blow up the errors significantly. But rarely do we see these cases.
Furthermore, there are various ways to interpret the data, and hence motivations for different interpolation schemes.
The \textit{error in the choice of interpolation scheme} will in nearly all cases be larger than the difference in precision by 
these two interpolation schemes.

We see that iteratively trying to decrease the error on the solution found by the LU decomposition is redundant. 
This is because we have employed implicit pivoting into our LU decomposition. This ensures that the maximum precision is obtained for
the datatype which is stored in the matrix. 



\subsection{Code for Question 2}
The following code was used for parts (2a)--(2d):
The imported code is given in  section \ref{appendix_section}.
\lstinputlisting[language=Python]{\CodeDir/Vandermonde_all_code.txt}


\section{Code in helper classes and functions}
\label{appendix_section}
The following code was imported in the exercises:
\lstinputlisting[language=Python]{\CodeDir/matrix.txt}
\lstinputlisting[language=Python]{\CodeDir/bisection.txt}

\end{document}

